\section{Espaces d’approximation sur les maillages \ddfv} \label{sec:unknowns}

La méthode \ddfv associe, à chaque maille primale plate $\primalplat{K} \in \ensprimalplat$, une inconnue scalaire $u_{\primalplat{K}}$ et, à chaque maille duale plate $\mailledualeplate{K} \in \ensdualplat$, une inconnue scalaire $u_{\mailledualeplate{K}}$. Ainsi l'ensemble des inconnues associé au maillage \ddfv $\ensddfvplat$ est défini de la manière suivante :
\begin{equation} \label{eq:ddfv_unknown}
u^\ensddfvplat \in \R^\ensddfvplat
\iff
u^\ensddfvplat \defeq \big( u^\ensprimalplat, u^\ensdualplat \big) 
\quad
\text{avec}
\quad
u^\ensprimalplat = (u_{\primalplat{K}})_{\primalplat{K} \in \ensprimalplat},
\quad
u^\ensdualplat = (u_{\mailledualeplate{K}})_{\mailledualeplate{K} \in \ensdualplat}.
\end{equation}
On a environ deux fois plus d'inconnues que dans la méthode VF4 mais cela permet de définir une approximation complète du gradient et pas uniquement dans la direction normale. Ceci rend la méthode beaucoup plus robuste. Les espaces d'approximation des fonctions discrètes $u^\ensprimalplat$ et $u^\ensdualplat$ dans le sens de la définition \ref{eq:ddfv_unknown} sont notés $\R^\ensprimalplat$ et $\R^\ensdualplat$ respectivement. En introduisant le projecteur
\begin{equation} \label{eq:projddfvvar}
\fonction[\projddfvvar]{\R^\ensddfvplat}{\L^2(\surflisse)}{u^\ensddfvplat}{\frac{1}{2} \sum_{\primal{K} \in \ensprimal} \indicatrice_{\primal{K}} \, u_{\primalplat{K}} + \frac{1}{2} \sum_{\mailleduale{K} \in \ensdual} \indicatrice_{\mailleduale{K}} \, u_{\mailledualeplate{K}}},
\end{equation}
on associe à la fonction discrète $u^\ensddfvplat$ la fonction constante par morceaux $\projddfvvar u^\ensddfvplat$. On utilise implicitement la correspondance entre une maille primale plate $\primalplat{K}$ et la maille primale courbe $\primal{K}$, et de même pour les mailles duales. Avec cette définition, nous utilisons simultanément deux reconstructions différentes basées soit sur les valeurs primales soit sur les valeurs duales. L'espace des solutions approchées introduit dans \cite{krell_developpement_2021} est dénoté par $\H_\ensddfvplat$ :
\[
\H_\ensddfvplat = \ens[\Big]{u \in \L^1(\surflisse) \tq \exists u^\ensddfvplat \in \R^\ensddfvplat \text{ tel que } u = \projddfvvar u^\ensddfvplat}.
\]
Nous pouvons ainsi définir la norme $\L^r$ de la fonction $\projddfvvar u^\ensddfvplat$ pour tout $r \in \intervalleff{1}{+\infty}$. 

On définit également la fonction vectorielle discrète $\xi^\ensdiamantplat \in (\R^2)^\ensdiamantplat$ par $\xi^\ensdiamantplat = (\xi^{\diamantplat{D}})_{\diamantplat{D} \in \ensdiamantplat}$, avec $\xi^{\diamantplat{D}} \in \R^2$. De la même manière que précédemment, en introduisant le projecteur
\begin{equation} \label{eq:projdiamantvar}
\fonction[\projdiamantvar]{(\R^2)^\ensdiamantplat}{(\L^2(\surflisse))^2}{\xi^\ensdiamantplat}{\sum_{\diamant{D} \in \ensdiamant} \indicatrice_{\diamant{D}} \, \xi^{\diamantplat{D}}},
\end{equation}
on associe à la fonction discrète $\xi^\ensdiamantplat$ la fonction constante par diamants courbes $\projdiamantvar \xi^\ensdiamantplat$. À nouveau, on utilise implicitement la correspondance entre un diamant plat $\diamantplat{D}$ et le diamant courbe $\diamant{D}$.